\Needspace{8\baselineskip}\originalchapter{箍缩、磁剪切与理想 kink}\label{ch:pinchstable}
\originalsection{简单 pinch 的压力驱动}
z pinch 无纵向场, 磁力线弯曲仅靠环向 B\ensuremath{\theta}, 容易产生轴对称 sausage 与 m=1侧向弯曲. 课程 Lecture17 在适当局部短波与边界条件下给 m=0 必要稳定式
\begin{equation}
 \frac{r p'}p\ge-\frac{4\gamma}{2+\gamma\beta(r)},\quad
 \beta(r)=\frac{2\mu_0p}{B_\theta^2},
\end{equation}
以及 m\ensuremath{\ne}0 的局部条件 $\dd(rB_\theta^2)/\dd r\le0$. 这些式由能量泛函的可允许局部位移最小化得来, 指定的硬芯与外壁必须另检查全局边界. 对沿 z 的场很强的 \ensuremath{\theta} pinch, 良曲率缺失但线张力可以稳定非轴对称扰动; 不可将无 Bz 的 z pinch 判据照搬.

\originalsection{螺旋 pinch 的共振面}
扰动 $\exp[\ii(m\theta+kz)]$, 沿场导数含 $F=mB_\theta/r+kB_z$. F=0 处弯曲恢复项减弱, \ensuremath{\xi}平行分量的形式解可能奇异. 在共振附近 x=r\ensuremath{-}rs, 能量约化出现 $x^2|\xi'|^2-D_s|\xi|^2$. Euler 方程的指数根在 $D_s>1/4$ 为复, 可构造局域负能位移. 换回平衡量得到 Suydam 必要局部稳定条件
\begin{equation}
 \frac{rB_z^2}{8\mu_0}\left(\frac{q'}q\right)^2+p'>0.
\end{equation}
剪切的正项抵消通常为负的压力梯度. Suydam 通过仅排除局域 interchange, 外 kink 或有限 m 模式仍可能不稳定.

\originalsection{安全因子和外 kink}
周期长度 $2\pi R_0$ 螺旋柱, $k=-n/R_0$, F\ensuremath{\propto}m\ensuremath{-}nq. 外 kink 将柱整体弯曲, 自由边界的真空场与导体壁提供额外弯曲能. 在简化圆柱模型 m=n=1的 Kruskal--Shafranov 约束说明 $q_{\rm edge}$ 低于约1容易失稳, 但真实电流剖面、边界、环形效应改变阈值. 内 kink 需共振面与内部剖面, 它不是外 kink 换个壁半径. 理解 q 的作用应结合模式整数、剖面和壁, 不是记住单个数字.
